\documentclass[aspectratio=169]{beamer}
\usepackage{amsmath}
\usepackage{amssymb}
\usepackage{tikz}
\usepackage{xcolor}
\usetikzlibrary{shapes,backgrounds,patterns}

\usepackage{pgfplots}
\pgfplotsset{compat=1.18, legend cell align=left}
\usetikzlibrary{arrows.meta,decorations.pathreplacing,calc,positioning}
\usepackage{booktabs}
\usepackage{pifont}
\usepgfplotslibrary{fillbetween}

% royalblue isn't a base xcolor name, so define it first (this is the default accent)
\definecolor{royalblue}{RGB}{65,105,225}

% ============================================================
%  ACCENT COLOR: change ONE line below to switch the theme
%  Choose one of: royalblue , red , black
% ============================================================
\colorlet{accent}{royalblue}
% \colorlet{accent}{red}
% \colorlet{accent}{black}
% ============================================================

% colorblind-safe, high-contrast second color (Okabe-Ito vermillion)
\definecolor{hl}{RGB}{213,94,0}

\usetheme{default}
\setbeamercolor{background canvas}{bg=white}
\setbeamercolor{normal text}{fg=black}
\setbeamercolor{structure}{fg=accent}
\setbeamercolor{frametitle}{fg=accent,bg=white}
\setbeamercolor{title}{fg=accent}
\setbeamercolor{section in toc}{fg=black}
\setbeamercolor{block title}{fg=black,bg=gray!10}
\setbeamercolor{block body}{fg=black,bg=gray!5}
\setbeamercolor{block title example}{fg=black,bg=green!10}
\setbeamercolor{block body example}{fg=black,bg=green!5}
\setbeamercolor{block title alerted}{fg=black,bg=red!10}
\setbeamercolor{block body alerted}{fg=black,bg=red!5}

\setbeamertemplate{navigation symbols}{}
\setbeamertemplate{footline}{
    \hfill\makebox[5pt]{%
        \scriptsize\insertframenumber}
    \hspace*{2em}
    \vspace{0.5cm}
}

\newcommand{\cbox}[2]{\colorbox{#1}{\ensuremath{\displaystyle #2}}}

\pgfplotsset{
  miniplot/.style={
    width=4.0cm, height=3.0cm,
    axis lines=middle, xtick=\empty, ytick=\empty,
    samples=120, thick, clip=true
  }
}

\definecolor{gr}{RGB}{0,158,115}


% curve and helpers for the three-rule picture
\pgfmathdeclarefunction{fc}{1}{\pgfmathparse{0.55+1.25*sin(deg(0.85*#1))}}
\pgfmathdeclarefunction{gc}{1}{\pgfmathparse{0.55+0.3*sin(deg(#1))+1.1*exp(-4*(#1-2.6)^2)}}
\pgfplotsset{
  rule2/.style={
    width=5.4cm, height=4.2cm, axis lines=middle, xmin=0, xmax=4.3, ymin=0, ymax=1.95,
    xtick=\empty, ytick=\empty, thick, clip=false,
    title style={font=\small}
  }
}
\pgfplotsset{
  rule/.style={
    width=5.4cm, height=4.2cm, axis lines=middle, xmin=0, xmax=4.3, ymin=0, ymax=1.95,
    xtick=\empty, ytick=\empty, thick, clip=false, samples=100,
    title style={font=\small}
  }
}

\title{Integration with SymPy and SciPy}
\subtitle{Exact and numerical integrals in one dimension}
\author{Sharareh Sayyad}
\institute{Washington State University}
\date{\today}

\begin{document}
\frame{\titlepage}

% ------------------------------------------------------------
\begin{frame}{Two routes to a definite integral}
\[ \int_a^b f(x)\,dx \;=\; F(b) - F(a), \qquad F'(x) = f(x) \]

\vspace{0.2em}
\begin{columns}[T]
\begin{column}{0.48\textwidth}
\begin{block}{Exact: SymPy}
\small
Find an antiderivative $F$.\\[2pt]
\footnotesize
\textcolor{gr}{\ding{51}} a formula, no rounding\\
\textcolor{hl}{\ding{55}} not every $f$ has one
\end{block}
\end{column}
\begin{column}{0.48\textwidth}
\begin{block}{Numerical: SciPy}
\small
Add up $f$ at sample points.\\[2pt]
\footnotesize
\textcolor{gr}{\ding{51}} always gives a number\\
\textcolor{hl}{\ding{55}} only approximate: check the error
\end{block}
\end{column}
\end{columns}

\vspace{0.8em}
\centering
\cbox{accent!12}{\text{Same idea as derivatives: \textbf{formula} vs.\ \textbf{data}.}}
\end{frame}

% ------------------------------------------------------------
\begin{frame}{Numerical integration: three rules}
\centering
\begin{tikzpicture}
% trapezoid
\begin{axis}[rule, title={Trapezoidal rule}, name=T]
\addplot[draw=gr, very thick, fill=gr!15, domain=0.6:3.6, samples=2]
  {fc(0.6) + (fc(3.6)-fc(0.6))/3*(x-0.6)} \closedcycle;
\addplot[black, domain=0.15:4.1] {fc(x)};
\addplot[only marks, mark=*, mark size=2pt] coordinates {(0.6,{fc(0.6)}) (3.6,{fc(3.6)})};
\node[below] at (axis cs:0.6,0) {\footnotesize $a$};
\node[below] at (axis cs:3.6,0) {\footnotesize $b$};
\end{axis}
% simpson
\begin{axis}[rule, title={Simpson's rule}, at={(T.east)}, anchor=west, xshift=0.6cm, name=S]
\addplot[draw=gr, very thick, fill=gr!15, domain=0.6:3.6]
  {fc(0.6)*(x-2.1)*(x-3.6)/((0.6-2.1)*(0.6-3.6))
 + fc(2.1)*(x-0.6)*(x-3.6)/((2.1-0.6)*(2.1-3.6))
 + fc(3.6)*(x-0.6)*(x-2.1)/((3.6-0.6)*(3.6-2.1))} \closedcycle;
\addplot[black, domain=0.15:4.1] {fc(x)};
\draw[dashed, thin] (axis cs:2.1,0) -- (axis cs:2.1,{fc(2.1)});
\addplot[only marks, mark=*, mark size=2pt] coordinates {(0.6,{fc(0.6)}) (2.1,{fc(2.1)}) (3.6,{fc(3.6)})};
\node[below] at (axis cs:0.6,0) {\footnotesize $a$};
\node[below] at (axis cs:2.1,0) {\footnotesize $x_1$};
\node[below] at (axis cs:3.6,0) {\footnotesize $b$};
\end{axis}
% gauss
\begin{axis}[rule, title={Gauss quadrature}, at={(S.east)}, anchor=west, xshift=0.6cm]
% cubic through the curve: 2-point Gauss is exact for cubics, so this area = (b-a)/2 (f(x1)+f(x2))
\addplot[draw=gr, very thick, fill=gr!15, domain=0.6:3.6]
  {fc(0.6)*(x-1.23397)*(x-2.96603)*(x-3.6)/((0.6-1.23397)*(0.6-2.96603)*(0.6-3.6))
 + fc(1.23397)*(x-0.6)*(x-2.96603)*(x-3.6)/((1.23397-0.6)*(1.23397-2.96603)*(1.23397-3.6))
 + fc(2.96603)*(x-0.6)*(x-1.23397)*(x-3.6)/((2.96603-0.6)*(2.96603-1.23397)*(2.96603-3.6))
 + fc(3.6)*(x-0.6)*(x-1.23397)*(x-2.96603)/((3.6-0.6)*(3.6-1.23397)*(3.6-2.96603))} \closedcycle;
\addplot[black, domain=0.15:4.1] {fc(x)};
\draw[dashed, thin] (axis cs:1.23397,0) -- (axis cs:1.23397,{fc(1.23397)});
\draw[dashed, thin] (axis cs:2.96603,0) -- (axis cs:2.96603,{fc(2.96603)});
\addplot[only marks, mark=*, mark size=2pt] coordinates {(1.23397,{fc(1.23397)}) (2.96603,{fc(2.96603)})};
\node[below] at (axis cs:0.6,0) {\footnotesize $a$};
\node[below] at (axis cs:1.23397,0) {\footnotesize $x_1$};
\node[below] at (axis cs:2.96603,0) {\footnotesize $x_2$};
\node[below] at (axis cs:3.6,0) {\footnotesize $b$};
\end{axis}
\end{tikzpicture}

\vspace{0.3em}
\footnotesize
\begin{columns}[T]
\begin{column}{0.31\textwidth}\centering
straight line through the ends
\end{column}
\begin{column}{0.31\textwidth}\centering
parabola through 3 points
\end{column}
\begin{column}{0.31\textwidth}\centering
best points $x_i$ and weights $w_i$\\
\scriptsize $x_i$: roots of Legendre $P_n$\\
($n=2$: $\pm 1/\sqrt{3}$ on $[-1,1]$)
\end{column}
\end{columns}

\vspace{0.6em}
\begin{columns}[T]
\begin{column}{0.31\textwidth}\centering
$E \approx -\dfrac{(b-a)\,h^2}{12}\, f''(\xi)$
\end{column}
\begin{column}{0.31\textwidth}\centering
$E \approx -\dfrac{(b-a)\,h^4}{180}\, f^{(4)}(\xi)$
\end{column}
\begin{column}{0.31\textwidth}\centering
$E = \dfrac{(b-a)^5}{4320}\, f^{(4)}(\xi)$ \ \ ($n=2$)
\end{column}
\end{columns}

\vspace{0.6em}
\centering
\cbox{accent!12}{\text{equal spacing } h = \tfrac{b-a}{n} \quad\text{vs.}\quad \text{Gauss chooses where to sample}}
\end{frame}

% ------------------------------------------------------------
\begin{frame}{Many points: split $[a,b]$ into pieces}
\centering
\begin{tikzpicture}
\begin{axis}[rule2, title={Trapezoid, $n=8$}, name=T2]
\addplot[draw=gr, very thick, fill=gr!15] coordinates {(0.4000,0.6668) (0.8250,0.7704) (1.2500,0.8354) (1.6750,0.8843) (2.1000,1.2136) (2.5250,1.7990) (2.9500,1.2810) (3.3750,0.5802) (3.8000,0.3699)} \closedcycle;
\draw[thin, gr!70!black] (axis cs:0.8250,0) -- (axis cs:0.8250,0.7704);
\draw[thin, gr!70!black] (axis cs:1.2500,0) -- (axis cs:1.2500,0.8354);
\draw[thin, gr!70!black] (axis cs:1.6750,0) -- (axis cs:1.6750,0.8843);
\draw[thin, gr!70!black] (axis cs:2.1000,0) -- (axis cs:2.1000,1.2136);
\draw[thin, gr!70!black] (axis cs:2.5250,0) -- (axis cs:2.5250,1.7990);
\draw[thin, gr!70!black] (axis cs:2.9500,0) -- (axis cs:2.9500,1.2810);
\draw[thin, gr!70!black] (axis cs:3.3750,0) -- (axis cs:3.3750,0.5802);
\addplot[black, domain=0.15:4.1, samples=150] {gc(x)};
\addplot[only marks, mark=*, mark size=1.4pt] coordinates {(0.4000,0.6668) (0.8250,0.7704) (1.2500,0.8354) (1.6750,0.8843) (2.1000,1.2136) (2.5250,1.7990) (2.9500,1.2810) (3.3750,0.5802) (3.8000,0.3699)};
\node[below] at (axis cs:0.4,0) {\footnotesize $a$};
\node[below] at (axis cs:3.8,0) {\footnotesize $b$};
\end{axis}
\begin{axis}[rule2, title={Simpson, $n=8$}, at={(T2.east)}, anchor=west, xshift=0.5cm, name=S2]
\addplot[draw=gr, very thick, fill=gr!15, domain=0.4000:1.2500, samples=30] {gc(0.4000)*(x-0.8250)*(x-1.2500)/((0.4000-0.8250)*(0.4000-1.2500))+gc(0.8250)*(x-0.4000)*(x-1.2500)/((0.8250-0.4000)*(0.8250-1.2500))+gc(1.2500)*(x-0.4000)*(x-0.8250)/((1.2500-0.4000)*(1.2500-0.8250))} \closedcycle;
\addplot[draw=gr, very thick, fill=gr!15, domain=1.2500:2.1000, samples=30] {gc(1.2500)*(x-1.6750)*(x-2.1000)/((1.2500-1.6750)*(1.2500-2.1000))+gc(1.6750)*(x-1.2500)*(x-2.1000)/((1.6750-1.2500)*(1.6750-2.1000))+gc(2.1000)*(x-1.2500)*(x-1.6750)/((2.1000-1.2500)*(2.1000-1.6750))} \closedcycle;
\addplot[draw=gr, very thick, fill=gr!15, domain=2.1000:2.9500, samples=30] {gc(2.1000)*(x-2.5250)*(x-2.9500)/((2.1000-2.5250)*(2.1000-2.9500))+gc(2.5250)*(x-2.1000)*(x-2.9500)/((2.5250-2.1000)*(2.5250-2.9500))+gc(2.9500)*(x-2.1000)*(x-2.5250)/((2.9500-2.1000)*(2.9500-2.5250))} \closedcycle;
\addplot[draw=gr, very thick, fill=gr!15, domain=2.9500:3.8000, samples=30] {gc(2.9500)*(x-3.3750)*(x-3.8000)/((2.9500-3.3750)*(2.9500-3.8000))+gc(3.3750)*(x-2.9500)*(x-3.8000)/((3.3750-2.9500)*(3.3750-3.8000))+gc(3.8000)*(x-2.9500)*(x-3.3750)/((3.8000-2.9500)*(3.8000-3.3750))} \closedcycle;
\draw[thin, gr!70!black] (axis cs:1.2500,0) -- (axis cs:1.2500,0.8354);
\draw[thin, gr!70!black] (axis cs:2.1000,0) -- (axis cs:2.1000,1.2136);
\draw[thin, gr!70!black] (axis cs:2.9500,0) -- (axis cs:2.9500,1.2810);
\draw[thin, dashed, gr!70!black] (axis cs:0.8250,0) -- (axis cs:0.8250,0.7704);
\draw[thin, dashed, gr!70!black] (axis cs:1.6750,0) -- (axis cs:1.6750,0.8843);
\draw[thin, dashed, gr!70!black] (axis cs:2.5250,0) -- (axis cs:2.5250,1.7990);
\draw[thin, dashed, gr!70!black] (axis cs:3.3750,0) -- (axis cs:3.3750,0.5802);
\addplot[black, domain=0.15:4.1, samples=150] {gc(x)};
\addplot[only marks, mark=*, mark size=1.4pt] coordinates {(0.4000,0.6668) (0.8250,0.7704) (1.2500,0.8354) (1.6750,0.8843) (2.1000,1.2136) (2.5250,1.7990) (2.9500,1.2810) (3.3750,0.5802) (3.8000,0.3699)};
\node[below] at (axis cs:0.4,0) {\footnotesize $a$};
\node[below] at (axis cs:3.8,0) {\footnotesize $b$};
\end{axis}
\begin{axis}[rule2, title={\texttt{quad}: adaptive Gauss}, at={(S2.east)}, anchor=west, xshift=0.5cm]
\addplot[draw=gr, very thick, fill=gr!15, domain=0.4000:0.8250, samples=30] {gc(0.4000)*(x-0.4898)*(x-0.7352)*(x-0.8250)/((0.4000-0.4898)*(0.4000-0.7352)*(0.4000-0.8250))+gc(0.4898)*(x-0.4000)*(x-0.7352)*(x-0.8250)/((0.4898-0.4000)*(0.4898-0.7352)*(0.4898-0.8250))+gc(0.7352)*(x-0.4000)*(x-0.4898)*(x-0.8250)/((0.7352-0.4000)*(0.7352-0.4898)*(0.7352-0.8250))+gc(0.8250)*(x-0.4000)*(x-0.4898)*(x-0.7352)/((0.8250-0.4000)*(0.8250-0.4898)*(0.8250-0.7352))} \closedcycle;
\addplot[draw=gr, very thick, fill=gr!15, domain=0.8250:1.2500, samples=30] {gc(0.8250)*(x-0.9148)*(x-1.1602)*(x-1.2500)/((0.8250-0.9148)*(0.8250-1.1602)*(0.8250-1.2500))+gc(0.9148)*(x-0.8250)*(x-1.1602)*(x-1.2500)/((0.9148-0.8250)*(0.9148-1.1602)*(0.9148-1.2500))+gc(1.1602)*(x-0.8250)*(x-0.9148)*(x-1.2500)/((1.1602-0.8250)*(1.1602-0.9148)*(1.1602-1.2500))+gc(1.2500)*(x-0.8250)*(x-0.9148)*(x-1.1602)/((1.2500-0.8250)*(1.2500-0.9148)*(1.2500-1.1602))} \closedcycle;
\addplot[draw=gr, very thick, fill=gr!15, domain=1.2500:1.6750, samples=30] {gc(1.2500)*(x-1.3398)*(x-1.5852)*(x-1.6750)/((1.2500-1.3398)*(1.2500-1.5852)*(1.2500-1.6750))+gc(1.3398)*(x-1.2500)*(x-1.5852)*(x-1.6750)/((1.3398-1.2500)*(1.3398-1.5852)*(1.3398-1.6750))+gc(1.5852)*(x-1.2500)*(x-1.3398)*(x-1.6750)/((1.5852-1.2500)*(1.5852-1.3398)*(1.5852-1.6750))+gc(1.6750)*(x-1.2500)*(x-1.3398)*(x-1.5852)/((1.6750-1.2500)*(1.6750-1.3398)*(1.6750-1.5852))} \closedcycle;
\addplot[draw=gr, very thick, fill=gr!15, domain=1.6750:2.1000, samples=30] {gc(1.6750)*(x-1.7648)*(x-2.0102)*(x-2.1000)/((1.6750-1.7648)*(1.6750-2.0102)*(1.6750-2.1000))+gc(1.7648)*(x-1.6750)*(x-2.0102)*(x-2.1000)/((1.7648-1.6750)*(1.7648-2.0102)*(1.7648-2.1000))+gc(2.0102)*(x-1.6750)*(x-1.7648)*(x-2.1000)/((2.0102-1.6750)*(2.0102-1.7648)*(2.0102-2.1000))+gc(2.1000)*(x-1.6750)*(x-1.7648)*(x-2.0102)/((2.1000-1.6750)*(2.1000-1.7648)*(2.1000-2.0102))} \closedcycle;
\addplot[draw=gr, very thick, fill=gr!15, domain=2.1000:2.3125, samples=30] {gc(2.1000)*(x-2.1449)*(x-2.2676)*(x-2.3125)/((2.1000-2.1449)*(2.1000-2.2676)*(2.1000-2.3125))+gc(2.1449)*(x-2.1000)*(x-2.2676)*(x-2.3125)/((2.1449-2.1000)*(2.1449-2.2676)*(2.1449-2.3125))+gc(2.2676)*(x-2.1000)*(x-2.1449)*(x-2.3125)/((2.2676-2.1000)*(2.2676-2.1449)*(2.2676-2.3125))+gc(2.3125)*(x-2.1000)*(x-2.1449)*(x-2.2676)/((2.3125-2.1000)*(2.3125-2.1449)*(2.3125-2.2676))} \closedcycle;
\addplot[draw=gr, very thick, fill=gr!15, domain=2.3125:2.5250, samples=30] {gc(2.3125)*(x-2.3574)*(x-2.4801)*(x-2.5250)/((2.3125-2.3574)*(2.3125-2.4801)*(2.3125-2.5250))+gc(2.3574)*(x-2.3125)*(x-2.4801)*(x-2.5250)/((2.3574-2.3125)*(2.3574-2.4801)*(2.3574-2.5250))+gc(2.4801)*(x-2.3125)*(x-2.3574)*(x-2.5250)/((2.4801-2.3125)*(2.4801-2.3574)*(2.4801-2.5250))+gc(2.5250)*(x-2.3125)*(x-2.3574)*(x-2.4801)/((2.5250-2.3125)*(2.5250-2.3574)*(2.5250-2.4801))} \closedcycle;
\addplot[draw=gr, very thick, fill=gr!15, domain=2.5250:2.7375, samples=30] {gc(2.5250)*(x-2.5699)*(x-2.6926)*(x-2.7375)/((2.5250-2.5699)*(2.5250-2.6926)*(2.5250-2.7375))+gc(2.5699)*(x-2.5250)*(x-2.6926)*(x-2.7375)/((2.5699-2.5250)*(2.5699-2.6926)*(2.5699-2.7375))+gc(2.6926)*(x-2.5250)*(x-2.5699)*(x-2.7375)/((2.6926-2.5250)*(2.6926-2.5699)*(2.6926-2.7375))+gc(2.7375)*(x-2.5250)*(x-2.5699)*(x-2.6926)/((2.7375-2.5250)*(2.7375-2.5699)*(2.7375-2.6926))} \closedcycle;
\addplot[draw=gr, very thick, fill=gr!15, domain=2.7375:2.9500, samples=30] {gc(2.7375)*(x-2.7824)*(x-2.9051)*(x-2.9500)/((2.7375-2.7824)*(2.7375-2.9051)*(2.7375-2.9500))+gc(2.7824)*(x-2.7375)*(x-2.9051)*(x-2.9500)/((2.7824-2.7375)*(2.7824-2.9051)*(2.7824-2.9500))+gc(2.9051)*(x-2.7375)*(x-2.7824)*(x-2.9500)/((2.9051-2.7375)*(2.9051-2.7824)*(2.9051-2.9500))+gc(2.9500)*(x-2.7375)*(x-2.7824)*(x-2.9051)/((2.9500-2.7375)*(2.9500-2.7824)*(2.9500-2.9051))} \closedcycle;
\addplot[draw=gr, very thick, fill=gr!15, domain=2.9500:3.3750, samples=30] {gc(2.9500)*(x-3.0398)*(x-3.2852)*(x-3.3750)/((2.9500-3.0398)*(2.9500-3.2852)*(2.9500-3.3750))+gc(3.0398)*(x-2.9500)*(x-3.2852)*(x-3.3750)/((3.0398-2.9500)*(3.0398-3.2852)*(3.0398-3.3750))+gc(3.2852)*(x-2.9500)*(x-3.0398)*(x-3.3750)/((3.2852-2.9500)*(3.2852-3.0398)*(3.2852-3.3750))+gc(3.3750)*(x-2.9500)*(x-3.0398)*(x-3.2852)/((3.3750-2.9500)*(3.3750-3.0398)*(3.3750-3.2852))} \closedcycle;
\addplot[draw=gr, very thick, fill=gr!15, domain=3.3750:3.8000, samples=30] {gc(3.3750)*(x-3.4648)*(x-3.7102)*(x-3.8000)/((3.3750-3.4648)*(3.3750-3.7102)*(3.3750-3.8000))+gc(3.4648)*(x-3.3750)*(x-3.7102)*(x-3.8000)/((3.4648-3.3750)*(3.4648-3.7102)*(3.4648-3.8000))+gc(3.7102)*(x-3.3750)*(x-3.4648)*(x-3.8000)/((3.7102-3.3750)*(3.7102-3.4648)*(3.7102-3.8000))+gc(3.8000)*(x-3.3750)*(x-3.4648)*(x-3.7102)/((3.8000-3.3750)*(3.8000-3.4648)*(3.8000-3.7102))} \closedcycle;
\draw[thin, gr!70!black] (axis cs:0.8250,0) -- (axis cs:0.8250,0.7704);
\draw[thin, gr!70!black] (axis cs:1.2500,0) -- (axis cs:1.2500,0.8354);
\draw[thin, gr!70!black] (axis cs:1.6750,0) -- (axis cs:1.6750,0.8843);
\draw[thin, gr!70!black] (axis cs:2.1000,0) -- (axis cs:2.1000,1.2136);
\draw[thin, gr!70!black] (axis cs:2.3125,0) -- (axis cs:2.3125,1.5615);
\draw[thin, gr!70!black] (axis cs:2.5250,0) -- (axis cs:2.5250,1.7990);
\draw[thin, gr!70!black] (axis cs:2.7375,0) -- (axis cs:2.7375,1.6878);
\draw[thin, gr!70!black] (axis cs:2.9500,0) -- (axis cs:2.9500,1.2810);
\draw[thin, gr!70!black] (axis cs:3.3750,0) -- (axis cs:3.3750,0.5802);
\addplot[black, domain=0.15:4.1, samples=150] {gc(x)};
\addplot[only marks, mark=*, mark size=1.4pt] coordinates {(0.4898,0.6911) (0.7352,0.7512) (0.9148,0.7877) (1.1602,0.8253) (1.3398,0.8439) (1.5852,0.8678) (1.7648,0.9119) (2.0102,1.0951) (2.1449,1.2823) (2.2676,1.4871) (2.3574,1.6311) (2.4801,1.7728) (2.5699,1.8083) (2.6926,1.7431) (2.7824,1.6184) (2.9051,1.3783) (3.0398,1.0879) (3.2852,0.6753) (3.4648,0.5099) (3.7102,0.3964)};
\node[below] at (axis cs:0.4,0) {\footnotesize $a$};
\node[below] at (axis cs:3.8,0) {\footnotesize $b$};
\end{axis}
\end{tikzpicture}

\vspace{0.3em}
\footnotesize
\begin{columns}[T]
\begin{column}{0.31\textwidth}\centering
a line on each piece
\end{column}
\begin{column}{0.31\textwidth}\centering
a parabola on each pair\\ so $n$ must be \textbf{even}
\end{column}
\begin{column}{0.31\textwidth}\centering
finer pieces near the peak
\end{column}
\end{columns}

\vspace{0.6em}
\begin{columns}[c]
\begin{column}{0.31\textwidth}\centering
$\dfrac{h}{2}\big(f_0 + 2f_1 + \dots + 2f_{n-1} + f_n\big)$
\end{column}
\begin{column}{0.31\textwidth}\centering
$\dfrac{h}{3}\big(f_0 + 4f_1 + 2f_2 + \dots + f_n\big)$
\end{column}
\begin{column}{0.31\textwidth}\centering
$\displaystyle\sum_{i=1}^{n} w_i\, f(x_i)$
\end{column}
\end{columns}
\end{frame}

\end{document}
